\documentclass[11pt,a4paper]{article}

\usepackage[utf8]{inputenc}
\usepackage[indonesian]{babel}
\usepackage[margin=2cm]{geometry}
\usepackage{amsmath,amssymb,amsfonts,amsthm}
\usepackage{booktabs}
\usepackage{longtable}
\usepackage{multirow}
\usepackage{enumitem}
\usepackage{titlesec}
\usepackage{float}
\usepackage{xcolor}
\usepackage{tcolorbox}
\usepackage{fancyhdr}
\usepackage{hyperref}
\usepackage{array}
\usepackage{listings}

\hypersetup{
    colorlinks=true,
    linkcolor=blue!70!black,
    citecolor=blue!70!black,
    urlcolor=blue!70!black,
    pdftitle={Panduan Lengkap Block Cipher: DES dan AES From Scratch},
    pdfauthor={MA4151 - Kriptografi},
}

\pagestyle{fancy}
\fancyhf{}
\rhead{\textbf{MA4151 - Kriptografi}}
\lhead{Panduan Teori \& Lembar Contoh Pengerjaan: DES \& AES}
\rfoot{Halaman \thepage}
\setlength{\headheight}{14pt}

% Colors definition
\definecolor{primaryblue}{RGB}{20, 50, 120}
\definecolor{accentblue}{RGB}{41, 128, 185}
\definecolor{boxbg}{RGB}{248, 250, 252}
\definecolor{darkgreen}{RGB}{27, 94, 32}
\definecolor{codebg}{RGB}{245, 245, 245}
\definecolor{probhead}{RGB}{180, 40, 40}

\tcbuselibrary{skins,breakable}

\tcbset{
    colback=boxbg,
    colframe=primaryblue,
    fonttitle=\bfseries,
    arc=2mm,
    left=4mm,
    right=4mm,
    top=3mm,
    bottom=3mm,
    breakable
}

\newtcolorbox{infobox}[1]{
    colback=blue!3!white,
    colframe=accentblue,
    fonttitle=\bfseries,
    title={#1},
    arc=1.5mm
}

\newtcolorbox{theorembox}[1]{
    colback=green!3!white,
    colframe=darkgreen,
    fonttitle=\bfseries,
    title={#1},
    arc=1.5mm
}

\newtcolorbox{problembox}[1]{
    colback=red!2!white,
    colframe=probhead,
    fonttitle=\bfseries,
    title={#1},
    arc=1.5mm
}

\lstset{
    backgroundcolor=\color{codebg},
    basicstyle=\ttfamily\small,
    breaklines=true,
    frame=single,
    rulecolor=\color{gray!40},
    keywordstyle=\color{primaryblue}\bfseries,
    commentstyle=\color{darkgreen}\itshape,
    stringstyle=\color{red!70!black}
}

\title{
    \vspace{-1cm}
    \textbf{\huge PANDUAN LENGKAP BLOCK CIPHER}\\[0.3cm]
    \textbf{\Large Data Encryption Standard (DES) \& Advanced Encryption Standard (AES)}\\[0.2cm]
    \large Teori, Tabel Standar Paten, dan Lembar Pengerjaan Soal Manual (Teks \& Desimal)
}
\author{\textbf{MA4151 -- Kriptografi}}
\date{\today}

\begin{document}

\maketitle

\begin{abstract}
Dokumen ini menyajikan panduan komprehensif mengenai kriptografi kunci simetris berbasis blok (\textit{block cipher}), berfokus pada dua standar utama: \textbf{Data Encryption Standard (DES)} dan \textbf{Advanced Encryption Standard (AES / Rijndael)}. Dokumen ini memuat seluruh \textbf{Tabel Paten Baku Standar NIST} ($IP$, $IP^{-1}$, $PC\text{-}1$, $PC\text{-}2$, $E$, $P$, $S_1 \dots S_8$), penjelasan representasi data berbasis \textbf{Karakter Teks Asli (ASCII)}, \textbf{Angka Desimal (0--255)}, dan \textbf{Bit Biner (0/1)}, serta lembar pengerjaan contoh soal langkah demi langkah bergaya tugas kuliah/PR.
\end{abstract}

\tableofcontents
\vspace{0.4cm}
\hrule
\vspace{0.4cm}

% =========================================================================
\section{Pengantar Sandi Blok \& Representasi Data}
% =========================================================================

Kriptografi kunci simetris (\textit{symmetric key cryptography}) memproses data menggunakan satu kunci rahasia yang sama untuk enkripsi dan dekripsi. Pada \textbf{Sandi Blok (Block Cipher)}, pesan dibagi menjadi blok-blok berukuran tetap:
\begin{itemize}[leftmargin=*]
    \item \textbf{DES:} Memproses blok berukuran \textbf{64 bit = 8 karakter teks (8 Byte)}.
    \item \textbf{AES:} Memproses blok berukuran \textbf{128 bit = 16 karakter teks (16 Byte)}.
\end{itemize}

\subsection{Bagaimana Karakter Teks Diubah Menjadi Angka \& Bit?}
Setiap karakter huruf atau simbol pada komputer memiliki kode angka \textbf{Desimal (ASCII)} antara $0$ sampai $255$, yang kemudian disimpan sebagai deretan \textbf{8 bit biner} ($0$ dan $1$):

\begin{table}[H]
\centering
\caption{Tabel Konversi Karakter Teks $\to$ Nilai Desimal $\to$ 8 Bit Biner}
\label{tab:ascii_primer}
\vspace{0.2cm}
\begin{tabular}{ccc}
\toprule
\textbf{Karakter Teks} & \textbf{Nilai Desimal (ASCII)} & \textbf{Representasi Biner (8 Bit)} \\
\midrule
\texttt{'K'} & 75  & \texttt{01001011} \\
\texttt{'O'} & 79  & \texttt{01001111} \\
\texttt{'M'} & 77  & \texttt{01001101} \\
\texttt{'P'} & 80  & \texttt{01010000} \\
\texttt{'U'} & 85  & \texttt{01010101} \\
\texttt{'T'} & 84  & \texttt{01010100} \\
\texttt{'E'} & 69  & \texttt{01000101} \\
\texttt{'R'} & 82  & \texttt{01010010} \\
\texttt{' '} (spasi) & 32 & \texttt{00100000} \\
\texttt{'!'} & 33  & \texttt{00100001} \\
\bottomrule
\end{tabular}
\end{table}

Dengan demikian, teks 8 karakter seperti \texttt{"KOMPUTER"} memiliki total $8 \times 8 = \mathbf{64\text{ bit biner}}$, yang siap diproses oleh algoritma DES. Teks 16 karakter seperti \texttt{"KRIPTOGRAFI ITB!"} memiliki total $16 \times 8 = \mathbf{128\text{ bit biner}}$, yang siap diproses oleh algoritma AES.

\subsection{Prinsip Claude Shannon: Kebingungan (Confusion) dan Difusi (Diffusion)}
\begin{itemize}[leftmargin=*]
    \item \textbf{Confusion (Kebingungan):} Membuat hubungan antara kunci rahasia dan teks tersandi menjadi rumit dan non-linear. Diwujudkan melalui \textbf{Kotak Substitusi (S-Box)}.
    \item \textbf{Diffusion (Difusi):} Menyebarkan pengaruh setiap bit teks terang ke seluruh blok ciphertext. Jika 1 bit teks terang berubah, rata-rata setengah dari seluruh bit ciphertext berubah (\textit{Avalanche Effect}). Diwujudkan melalui permutasi bit (\textbf{P-Box}), pergeseran baris (\textbf{ShiftRows}), dan pencampuran kolom (\textbf{MixColumns}).
\end{itemize}

\newpage
% =========================================================================
\section{Teori Data Encryption Standard (DES) \& Tabel Standar Paten}
% =========================================================================

DES memproses satu blok 64-bit (8 huruf) menggunakan kunci 64-bit (56 bit efektif + 8 bit paritas) melalui 16 putaran Jaringan Feistel.

\begin{infobox}{Apakah Tabel-Tabel pada DES Boleh Diacak Sendiri?}
\textbf{TIDAK BOLEH.} Seluruh tabel dalam DES ($IP$, $IP^{-1}$, $PC\text{-}1$, $PC\text{-}2$, $E$, $P$, dan $S_1 \dots S_8$) adalah **standar paten resmi internasional** (NIST FIPS PUB 46-3). Semua sistem di dunia wajib menggunakan tabel yang persis sama agar data yang dienkripsi dapat didekripsi dengan benar oleh sistem lain.
\end{infobox}

\subsection{Tabel Permutasi Awal ($IP$) \& Invers Permutasi Awal ($IP^{-1}$)}
Tabel $IP$ menentukan urutan baru posisi ke-64 bit teks terang. 
\begin{itemize}[leftmargin=*]
    \item \textbf{4 Baris Pertama (32 bit pertama)} $\to$ Membentuk register kiri $\mathbf{L_0}$ (semuanya adalah nomor bit genap).
    \item \textbf{4 Baris Kedua (32 bit kedua)} $\to$ Membentuk register kanan $\mathbf{R_0}$ (semuanya adalah nomor bit ganjil).
\end{itemize}

\begin{table}[H]
\centering
\caption{Tabel Baku Initial Permutation ($IP$) Standar FIPS 46-3 (64 bit $\to$ 64 bit)}
\label{tab:ip_table}
\vspace{0.2cm}
\begin{tabular}{|cccccccc|l}
\cline{1-8}
\textbf{58} & \textbf{50} & \textbf{42} & \textbf{34} & \textbf{26} & \textbf{18} & \textbf{10} & \textbf{2} & \multirow{4}{*}{$\Big\}\ \mathbf{32\text{ bit pertama}} \to \mathbf{L_0}$} \\
\textbf{60} & \textbf{52} & \textbf{44} & \textbf{36} & \textbf{28} & \textbf{20} & \textbf{12} & \textbf{4} & \\
\textbf{62} & \textbf{54} & \textbf{46} & \textbf{38} & \textbf{30} & \textbf{22} & \textbf{14} & \textbf{6} & \\
\textbf{64} & \textbf{56} & \textbf{48} & \textbf{40} & \textbf{32} & \textbf{24} & \textbf{16} & \textbf{8} & \\
\cline{1-8}
57 & 49 & 41 & 33 & 25 & 17 & 9 & 1 & \multirow{4}{*}{$\Big\}\ \mathbf{32\text{ bit kedua}} \to \mathbf{R_0}$} \\
59 & 51 & 43 & 35 & 27 & 19 & 11 & 3 & \\
61 & 53 & 45 & 37 & 29 & 21 & 13 & 5 & \\
63 & 55 & 47 & 39 & 31 & 23 & 15 & 7 & \\
\cline{1-8}
\end{tabular}
\end{table}

\noindent
\textbf{Cara membaca tabel di atas:}
\begin{itemize}[leftmargin=*]
    \item Bit ke-1 hasil permutasi $IP$ diambil dari **bit ke-58** teks asli.
    \item Bit ke-2 hasil permutasi $IP$ diambil dari **bit ke-50** teks asli.
    \item $\dots$ dan seterusnya sampai bit ke-64 diambil dari **bit ke-7** teks asli.
\end{itemize}

Di akhir Ronde 16, hasil permutasi dibalikkan menggunakan tabel pasangannya, yaitu $IP^{-1}$:

\begin{table}[H]
\centering
\caption{Tabel Baku Inverse Initial Permutation ($IP^{-1}$) Standar FIPS 46-3 (64 bit $\to$ 64 bit)}
\label{tab:ip_inv_table}
\vspace{0.2cm}
\begin{tabular}{|cccccccc|}
\hline
40 & 8 & 48 & 16 & 56 & 24 & 64 & 32 \\
39 & 7 & 47 & 15 & 55 & 23 & 63 & 31 \\
38 & 6 & 46 & 14 & 54 & 22 & 62 & 30 \\
37 & 5 & 45 & 13 & 53 & 21 & 61 & 29 \\
36 & 4 & 44 & 12 & 52 & 20 & 60 & 28 \\
35 & 3 & 43 & 11 & 51 & 19 & 59 & 27 \\
34 & 2 & 42 & 10 & 50 & 18 & 58 & 26 \\
33 & 1 & 41 & 9  & 49 & 17 & 57 & 25 \\
\hline
\end{tabular}
\end{table}

\subsection{Tabel Pembangkitan Kunci ($PC\text{-}1$, Shifts, $PC\text{-}2$)}

\begin{table}[H]
\centering
\caption{Tabel $PC\text{-}1$ (Membuang bit paritas 8, 16, 24, 32, 40, 48, 56, 64) $\to$ 56 bit ($C_0, D_0$)}
\vspace{0.2cm}
\begin{tabular}{|ccccccc|l}
\cline{1-7}
57 & 49 & 41 & 33 & 25 & 17 & 9  & \multirow{4}{*}{$\Big\}\ \mathbf{C_0\text{ (28 bit)}}$} \\
1  & 58 & 50 & 42 & 34 & 26 & 18 & \\
10 & 2  & 59 & 51 & 43 & 35 & 27 & \\
19 & 11 & 3  & 60 & 52 & 44 & 36 & \\
\cline{1-7}
63 & 55 & 47 & 39 & 31 & 23 & 15 & \multirow{4}{*}{$\Big\}\ \mathbf{D_0\text{ (28 bit)}}$} \\
7  & 62 & 54 & 46 & 38 & 30 & 22 & \\
14 & 6  & 61 & 53 & 45 & 37 & 29 & \\
21 & 13 & 5  & 28 & 20 & 12 & 4  & \\
\cline{1-7}
\end{tabular}
\end{table}

\begin{table}[H]
\centering
\caption{Jumlah Pergeseran Kiri Sirkular per Ronde \& Tabel Kompresi $PC\text{-}2$ (56 bit $\to$ 48 bit subkunci)}
\vspace{0.2cm}
\begin{minipage}{0.45\textwidth}
\centering
\textbf{Tabel Pergeseran (Shifts):}
\begin{tabular}{cc}
\toprule
\textbf{Ronde} & \textbf{Jumlah Pergeseran} \\
\midrule
1, 2, 9, 16 & 1 bit ke kiri \\
Lainnya (12 ronde) & 2 bit ke kiri \\
\bottomrule
\end{tabular}
\end{minipage}
\hfill
\begin{minipage}{0.5\textwidth}
\centering
\textbf{Tabel $PC\text{-}2$ (48 bit):}
\begin{tabular}{|cccccc|}
\hline
14 & 17 & 11 & 24 & 1  & 5  \\
3  & 28 & 15 & 6  & 21 & 10 \\
23 & 19 & 12 & 4  & 26 & 8  \\
16 & 7  & 27 & 20 & 13 & 2  \\
41 & 52 & 31 & 37 & 47 & 55 \\
30 & 40 & 51 & 45 & 33 & 48 \\
44 & 49 & 39 & 56 & 34 & 53 \\
46 & 42 & 50 & 36 & 29 & 32 \\
\hline
\end{tabular}
\end{minipage}
\end{table}

\subsection{Tabel Ekspansi $E$ \& Permutasi $P$}

\begin{table}[H]
\centering
\caption{Tabel Ekspansi $E$ (32 bit $\to$ 48 bit) \& Tabel Permutasi $P$ (32 bit $\to$ 32 bit)}
\vspace{0.2cm}
\begin{minipage}{0.48\textwidth}
\centering
\textbf{Tabel Ekspansi $E$:}
\begin{tabular}{|cccccc|}
\hline
32 & 1  & 2  & 3  & 4  & 5  \\
4  & 5  & 6  & 7  & 8  & 9  \\
8  & 9  & 10 & 11 & 12 & 13 \\
12 & 13 & 14 & 15 & 16 & 17 \\
16 & 17 & 18 & 19 & 20 & 21 \\
20 & 21 & 22 & 23 & 24 & 25 \\
24 & 25 & 26 & 27 & 28 & 29 \\
28 & 29 & 30 & 31 & 32 & 1  \\
\hline
\end{tabular}
\end{minipage}
\hfill
\begin{minipage}{0.48\textwidth}
\centering
\textbf{Tabel Permutasi $P$:}
\begin{tabular}{|cccccccc|}
\hline
16 & 7  & 20 & 21 & 29 & 12 & 28 & 17 \\
1  & 15 & 23 & 26 & 5  & 18 & 31 & 10 \\
2  & 8  & 24 & 14 & 32 & 27 & 3  & 9  \\
19 & 13 & 30 & 6  & 22 & 11 & 4  & 25 \\
\hline
\end{tabular}
\end{minipage}
\end{table}

\newpage
% =========================================================================
\section{Lembar Pengerjaan Contoh Soal Manual DES (Teks \& Desimal)}
% =========================================================================

\begin{problembox}{Soal Perhitungan Manual Enkripsi DES}
\textbf{Diketahui:}
\begin{itemize}[leftmargin=*]
    \item Teks Terang (\textit{Plaintext}): \texttt{"KOMPUTER"} (8 karakter teks)
    \item Kunci Rahasia (\textit{Key}): \texttt{"RAHASIA1"} (8 karakter teks)
\end{itemize}
\textbf{Tugas:}
\begin{enumerate}[leftmargin=*]
    \item Ubah teks terang dan kunci menjadi nilai desimal ASCII dan 64 bit biner.
    \item Lakukan Permutasi Awal ($IP$) pada teks terang untuk memperoleh $L_0$ dan $R_0$.
    \item Hitung subkunci Ronde 1 ($K_1$) menggunakan $PC\text{-}1$, pergeseran kiri 1 bit, dan $PC\text{-}2$.
    \item Hitung fungsi Feistel $f(R_0, K_1)$ dengan ekspansi $E$, XOR $K_1$, evaluasi 8 Kotak-S, dan permutasi $P$.
    \item Tentukan output Ronde 1 $(L_1, R_1)$ dan sajikan rekapitulasi seluruh 16 ronde hingga ciphertext akhir.
\end{enumerate}
\end{problembox}

\subsection*{Penyelesaian Langkah demi Langkah:}

\subsubsection*{Langkah 1: Konversi Teks Terang \& Kunci ke Biner}
\begin{itemize}[leftmargin=*]
    \item \textbf{Teks Terang \texttt{"KOMPUTER"}:}
    \begin{itemize}
        \item Desimal ASCII: $[75, 79, 77, 80, 85, 84, 69, 82]$
        \item Biner 64 bit:\\
        \texttt{01001011 01001111 01001101 01010000 01010101 01010100 01000101 01010010}
    \end{itemize}
    \item \textbf{Kunci \texttt{"RAHASIA1"}:}
    \begin{itemize}
        \item Desimal ASCII: $[82, 65, 72, 65, 83, 73, 65, 49]$
        \item Biner 64 bit:\\
        \texttt{01010010 01000001 01001000 01000001 01010011 01001001 01000001 00110001}
    \end{itemize}
\end{itemize}

\subsubsection*{Langkah 2: Permutasi Awal ($IP$) Teks Terang}
Petakan 64 bit teks terang dengan tabel $IP$ baku, lalu bagi dua menjadi 32 bit kiri dan 32 bit kanan:
\begin{align*}
\mathbf{L_0} &= \texttt{11111111 10111000 01110110 01010111}_2 \quad (\text{Desimal: } [255, 184, 118, 87]) \\
\mathbf{R_0} &= \texttt{00000000 00000000 00000111 10000011}_2 \quad (\text{Desimal: } [0, 0, 7, 131])
\end{align*}

\subsubsection*{Langkah 3: Pembangkitan Subkunci Ronde 1 ($K_1$)}
\begin{enumerate}[leftmargin=*]
    \item Terapkan $PC\text{-}1$ pada kunci 64-bit untuk memperoleh register 28-bit $C_0$ dan $D_0$:
    \begin{align*}
    C_0 &= \texttt{1000100 0111001 0110011 1100000}_2 \quad (28\text{ bit}) \\
    D_0 &= \texttt{0101000 1100011 0011111 0101001}_2 \quad (28\text{ bit})
    \end{align*}
    \item Geser ke kiri 1 bit ($\text{LS}_1$):
    \begin{align*}
    C_1 &= \texttt{0001000 1110010 1100111 1000001}_2 \\
    D_1 &= \texttt{1010001 1000110 0111110 1010010}_2
    \end{align*}
    \item Terapkan $PC\text{-}2$ pada $C_1 \parallel D_1$ untuk memperoleh subkunci 48-bit $K_1$:
    \begin{equation*}
    \mathbf{K_1} = \texttt{000010 101101 011100 100000 001011 001001 100001 100111}_2 \quad (48\text{ bit})
    \end{equation*}
\end{enumerate}

\subsubsection*{Langkah 4: Perhitungan Fungsi Feistel $f(R_0, K_1)$}
\begin{enumerate}[leftmargin=*]
    \item \textbf{Ekspansi $E(R_0)$:} Memperluas 32 bit $R_0$ menjadi 48 bit:
    \begin{equation*}
    E(R_0) = \texttt{100000 000000 000000 000000 000000 001111 110000 000010}_2
    \end{equation*}
    \item \textbf{XOR dengan $K_1$:}
    \begin{align*}
    A &= E(R_0) \oplus K_1 \\
    &= \texttt{100010 101101 011100 100000 001011 000110 010001 100101}_2 \quad (48\text{ bit})
    \end{align*}
    \item \textbf{Evaluasi 8 Kotak-S ($S_1$ s.d $S_8$):}
    \begin{table}[H]
    \centering
    \small
    \begin{tabular}{ccccccc}
    \toprule
    \textbf{Kotak-S} & \textbf{Bit Input (6-bit)} & \textbf{Bit Luar ($b_1 b_6$)} & \textbf{Baris Desimal} & \textbf{Bit Tengah ($b_2 b_3 b_4 b_5$)} & \textbf{Kolom Desimal} & \textbf{Output (4-bit)} \\
    \midrule
    $S_1$ & \texttt{100010} & \texttt{10} & 2 & \texttt{0001} & 1  & $1 = \texttt{0001}_2$ \\
    $S_2$ & \texttt{101101} & \texttt{11} & 3 & \texttt{0110} & 6  & $4 = \texttt{0100}_2$ \\
    $S_3$ & \texttt{011100} & \texttt{00} & 0 & \texttt{1110} & 14 & $0 = \texttt{0000}_2$ \\
    $S_4$ & \texttt{100000} & \texttt{10} & 2 & \texttt{0000} & 0  & $10 = \texttt{1010}_2$ \\
    $S_5$ & \texttt{001011} & \texttt{01} & 1 & \texttt{0101} & 5  & $0 = \texttt{0000}_2$ \\
    $S_6$ & \texttt{000110} & \texttt{00} & 0 & \texttt{0011} & 3  & $15 = \texttt{1111}_2$ \\
    $S_7$ & \texttt{010001} & \texttt{01} & 1 & \texttt{1000} & 8  & $14 = \texttt{1110}_2$ \\
    $S_8$ & \texttt{100101} & \texttt{11} & 3 & \texttt{0010} & 2  & $14 = \texttt{1110}_2$ \\
    \bottomrule
    \end{tabular}
    \end{table}
    Gabungkan ke-8 output 4-bit menjadi 32 bit:
    \begin{equation*}
    \text{Output Kotak-S} = \texttt{0001 0100 0000 1010 0000 1111 1110 1110}_2
    \end{equation*}
    \item \textbf{Permutasi $P$ (P-Box):}
    \begin{equation*}
    \mathbf{f(R_0, K_1)} = P(\text{Output Kotak-S}) = \texttt{11000010 11001101 01101010 01011110}_2
    \end{equation*}
\end{enumerate}

\subsubsection*{Langkah 5: Output Ronde 1 \& Rekapitulasi 16 Ronde}
\begin{align*}
\mathbf{L_1} &= R_0 = \texttt{00000000 00000000 00000111 10000011}_2 \quad (\text{Desimal: } [0, 0, 7, 131]) \\
\mathbf{R_1} &= L_0 \oplus f(R_0, K_1) = \texttt{00111101 01110101 00011100 00001001}_2 \quad (\text{Desimal: } [61, 117, 28, 9])
\end{align*}

\begin{table}[H]
\centering
\small
\caption{Rekapitulasi Nilai State Desimal $L_i$ dan $R_i$ Selama 16 Ronde DES}
\vspace{0.2cm}
\begin{tabular}{ccc}
\toprule
\textbf{Ronde ($i$)} & $\mathbf{L_i}$ \textbf{(Nilai 4 Byte Desimal)} & $\mathbf{R_i}$ \textbf{(Nilai 4 Byte Desimal)} \\
\midrule
0 (Awal) & $[255, 184, 118, 87]$ & $[0, 0, 7, 131]$ \\
1  & $[0, 0, 7, 131]$      & $[61, 117, 28, 9]$ \\
2  & $[61, 117, 28, 9]$    & $[73, 32, 63, 209]$ \\
3  & $[73, 32, 63, 209]$   & $[36, 183, 113, 161]$ \\
4  & $[36, 183, 113, 161]$ & $[23, 7, 71, 126]$ \\
5  & $[23, 7, 71, 126]$    & $[106, 248, 150, 72]$ \\
6  & $[106, 248, 150, 72]$ & $[106, 241, 46, 154]$ \\
7  & $[106, 241, 46, 154]$ & $[36, 34, 167, 68]$ \\
8  & $[36, 34, 167, 68]$   & $[194, 202, 237, 87]$ \\
9  & $[194, 202, 237, 87]$ & $[210, 85, 151, 4]$ \\
10 & $[210, 85, 151, 4]$   & $[43, 61, 125, 138]$ \\
11 & $[43, 61, 125, 138]$  & $[231, 94, 239, 140]$ \\
12 & $[231, 94, 239, 140]$ & $[118, 111, 59, 235]$ \\
13 & $[118, 111, 59, 235]$ & $[9, 145, 26, 170]$ \\
14 & $[9, 145, 26, 170]$   & $[63, 50, 207, 29]$ \\
15 & $[63, 50, 207, 29]$   & $[214, 147, 36, 149]$ \\
16 & $[214, 147, 36, 149]$ & $[28, 75, 115, 197]$ \\
\bottomrule
\end{tabular}
\end{table}

\textbf{Pertukaran Akhir 32-bit:} Gabungkan $(R_{16} \parallel L_{16})$.\\
\textbf{Teks Tersandi Akhir ($C = IP^{-1}(R_{16} \parallel L_{16})$):}
\begin{itemize}[leftmargin=*]
    \item Nilai Desimal (8 Byte): $\mathbf{[55, 180, 203, 80, 230, 12, 149, 163]}$
    \item Biner (64 bit): \texttt{00110111 10110100 11001011 01010000 11100110 00001100 10010101 10100011}
\end{itemize}

\newpage
% =========================================================================
\section{Teori Advanced Encryption Standard (AES / Rijndael)}
% =========================================================================

AES beroperasi pada blok \textbf{128 bit (16 karakter teks / 16 Byte)} yang disusun dalam matriks $4 \times 4$ byte (\textit{State Matrix}) berurutan kolom demi kolom (\textit{column-major}).

\subsection{Matriks State $4 \times 4$}
16 karakter pesan diposisikan sebagai berikut:
\begin{equation}
\text{State Karakter} = \begin{bmatrix}
\text{Karakter 0} & \text{Karakter 4} & \text{Karakter 8}  & \text{Karakter 12} \\
\text{Karakter 1} & \text{Karakter 5} & \text{Karakter 9}  & \text{Karakter 13} \\
\text{Karakter 2} & \text{Karakter 6} & \text{Karakter 10} & \text{Karakter 14} \\
\text{Karakter 3} & \text{Karakter 7} & \text{Karakter 11} & \text{Karakter 15}
\end{bmatrix}
\end{equation}

\subsection{Empat Transformasi Ronde AES}
\begin{enumerate}[leftmargin=*]
    \item \textbf{SubBytes:} Mengganti setiap byte dengan nilai tabel Kotak-S Rijndael (berdasarkan invers $GF(2^8)$ dan transformasi afina).
    \item \textbf{ShiftRows:} Menggeser baris secara melingkar ke kiri (Baris 0 geser 0, Baris 1 geser 1, Baris 2 geser 2, Baris 3 geser 3).
    \item \textbf{MixColumns:} Mengalikan tiap kolom dengan matriks konstanta MDS di Lapangan Galois $GF(2^8)$.
    \item \textbf{AddRoundKey:} Melakukan operasi XOR antara matriks State dengan matriks subkunci ronde ($4 \times 4$).
\end{enumerate}

\newpage
% =========================================================================
\section{Lembar Pengerjaan Contoh Soal Manual AES-128 (Teks \& Desimal)}
% =========================================================================

\begin{problembox}{Soal Perhitungan Manual Enkripsi AES-128}
\textbf{Diketahui:}
\begin{itemize}[leftmargin=*]
    \item Teks Terang (\textit{Plaintext}): \texttt{"KRIPTOGRAFI ITB!"} (16 karakter teks)
    \item Kunci Rahasia (\textit{Key}): \texttt{"KUNCI RAHASIAKU!"} (16 karakter teks)
\end{itemize}
\textbf{Tugas:}
\begin{enumerate}[leftmargin=*]
    \item Susun matriks State awal dan matriks subkunci $K_0$ dalam karakter teks dan nilai desimal ($0 \dots 255$).
    \item Hitung pembangkitan subkunci ronde 1 ($K_1$) dalam bentuk desimal.
    \item Hitung Initial AddRoundKey ($\text{State} \oplus K_0$).
    \item Hitung seluruh 4 transformasi pada Ronde 1: $\mathtt{SubBytes}$, $\mathtt{ShiftRows}$, $\mathtt{MixColumns}$ (tampilkan matriks desimal), dan $\mathtt{AddRoundKey}$.
    \item Tampilkan tabel rekapitulasi state akhir tiap ronde dan ciphertext akhir dalam desimal dan biner.
\end{enumerate}
\end{problembox}

\subsection*{Penyelesaian Langkah demi Langkah:}

\subsubsection*{Langkah 1: Matriks State Awal \& Subkunci $K_0$}
Teks terang 16 karakter: \texttt{'K', 'R', 'I', 'P', 'T', 'O', 'G', 'R', 'A', 'F', 'I', ' ', 'I', 'T', 'B', '!'} \\
Nilai Desimal ASCII: $[75, 82, 73, 80, 84, 79, 71, 82, 65, 70, 73, 32, 73, 84, 66, 33]$

Disusun dalam matriks $4 \times 4$ (\textit{column-major}):
\begin{equation*}
\text{State Awal (Karakter)} = \begin{bmatrix}
\text{'K'} & \text{'T'} & \text{'A'} & \text{'I'} \\
\text{'R'} & \text{'O'} & \text{'F'} & \text{'T'} \\
\text{'I'} & \text{'G'} & \text{'I'} & \text{'B'} \\
\text{'P'} & \text{'R'} & \text{' '} & \text{'!'}
\end{bmatrix}, \quad
\text{State Awal (Desimal)} = \begin{bmatrix}
75 & 84 & 65 & 73 \\
82 & 79 & 70 & 84 \\
73 & 71 & 73 & 66 \\
80 & 82 & 32 & 33
\end{bmatrix}
\end{equation*}

Kunci rahasia 16 karakter: \texttt{'K', 'U', 'N', 'C', 'I', ' ', 'R', 'A', 'H', 'A', 'S', 'I', 'A', 'K', 'U', '!'} \\
Nilai Desimal ASCII: $[75, 85, 78, 67, 73, 32, 82, 65, 72, 65, 83, 73, 65, 75, 85, 33]$
\begin{equation*}
\mathbf{K_0\text{ (Desimal)}} = \begin{bmatrix}
75 & 73 & 72 & 65 \\
85 & 32 & 65 & 75 \\
78 & 82 & 83 & 85 \\
67 & 65 & 73 & 33
\end{bmatrix}
\end{equation*}

\subsubsection*{Langkah 2: Pembangkitan Subkunci Ronde 1 ($K_1$)}
Melalui algoritma Key Expansion dengan $\mathtt{RotWord}$, $\mathtt{SubWord}$ (S-box), dan $\mathtt{Rcon}[1]$, diperoleh matriks subkunci ronde 1 dalam bentuk desimal:
\begin{equation*}
\mathbf{K_1\text{ (Desimal)}} = \begin{bmatrix}
249 & 176 & 248 & 185 \\
169 & 137 & 200 & 131 \\
179 & 225 & 178 & 231 \\
192 & 129 & 200 & 233
\end{bmatrix}
\end{equation*}

\subsubsection*{Langkah 3: Initial AddRoundKey ($\text{State} \oplus K_0$)}
Lakukan operasi XOR bitwise antara setiap elemen State Awal dengan $K_0$:
\begin{equation*}
\text{State}_0 = \begin{bmatrix}
75 \oplus 75 & 84 \oplus 73 & 65 \oplus 72 & 73 \oplus 65 \\
82 \oplus 85 & 79 \oplus 32 & 70 \oplus 65 & 84 \oplus 75 \\
73 \oplus 78 & 71 \oplus 82 & 73 \oplus 83 & 66 \oplus 85 \\
80 \oplus 67 & 82 \oplus 65 & 32 \oplus 73 & 33 \oplus 33
\end{bmatrix}
= \mathbf{\begin{bmatrix}
0  & 29  & 9   & 8 \\
7  & 111 & 7   & 31 \\
7  & 21  & 26  & 23 \\
19 & 19  & 105 & 0
\end{bmatrix}}
\end{equation*}

\subsubsection*{Langkah 4: Empat Transformasi Ronde 1}

\paragraph{4a. SubBytes:} Setiap nilai digantikan dengan nilai S-box (dalam desimal):
\begin{equation*}
\text{State}_{\text{SubBytes}} = \mathbf{\begin{bmatrix}
99  & 164 & 1   & 48 \\
197 & 168 & 197 & 192 \\
197 & 89  & 162 & 240 \\
125 & 125 & 249 & 99
\end{bmatrix}}
\end{equation*}

\paragraph{4b. ShiftRows:} Geser baris ke kiri secara melingkar (baris $0, 1, 2, 3$ geser $0, 1, 2, 3$ posisi):
\begin{equation*}
\text{State}_{\text{ShiftRows}} = \mathbf{\begin{bmatrix}
99  & 164 & 1   & 48 \\
168 & 197 & 192 & 197 \\
162 & 240 & 197 & 89 \\
99  & 125 & 125 & 249
\end{bmatrix}}
\end{equation*}

\paragraph{4c. MixColumns:} Perkalian dengan matriks MDS atas Lapangan Galois $GF(2^8)$:
\begin{equation*}
\text{State}_{\text{MixColumns}} = \mathbf{\begin{bmatrix}
228 & 138 & 225 & 148 \\
182 & 67  & 179 & 179 \\
49  & 29  & 215 & 87 \\
105 & 56  & 252 & 37
\end{bmatrix}}
\end{equation*}

\paragraph{4d. AddRoundKey ($K_1$):} Lakukan XOR dengan matriks subkunci $K_1$:
\begin{equation*}
\text{State}_{\text{Ronde 1 Akhir}} = \begin{bmatrix}
228 \oplus 249 & 138 \oplus 176 & 225 \oplus 248 & 148 \oplus 185 \\
182 \oplus 169 & 67 \oplus 137  & 179 \oplus 200 & 179 \oplus 131 \\
49 \oplus 179  & 29 \oplus 225  & 215 \oplus 178 & 87 \oplus 231 \\
105 \oplus 192 & 56 \oplus 129  & 252 \oplus 200 & 37 \oplus 233
\end{bmatrix}
= \mathbf{\begin{bmatrix}
29  & 58  & 25  & 45 \\
31  & 202 & 123 & 48 \\
130 & 252 & 101 & 176 \\
169 & 185 & 52  & 204
\end{bmatrix}}
\end{equation*}

\subsubsection*{Langkah 5: Rekapitulasi State Akhir Tiap Ronde \& Ciphertext Akhir}
\begin{table}[H]
\centering
\small
\caption{Rekapitulasi 16 Nilai Byte Desimal State Akhir Tiap Ronde AES-128}
\vspace{0.2cm}
\begin{tabular}{cl}
\toprule
\textbf{Ronde} & \textbf{Nilai 16 Byte Desimal State Akhir (Column-Major)} \\
\midrule
Awal (ARK 0) & $[0, 7, 7, 19, 29, 111, 21, 19, 9, 7, 26, 105, 8, 31, 23, 0]$ \\
Ronde 1  & $[29, 31, 130, 169, 58, 202, 252, 185, 25, 123, 101, 52, 45, 48, 176, 204]$ \\
Ronde 2  & $[188, 111, 241, 195, 201, 185, 151, 193, 247, 83, 172, 165, 80, 249, 128, 92]$ \\
Ronde 3  & $[160, 204, 188, 115, 59, 102, 180, 165, 174, 91, 140, 157, 173, 40, 199, 127]$ \\
Ronde 4  & $[201, 107, 143, 73, 239, 158, 98, 124, 76, 171, 158, 195, 25, 215, 60, 112]$ \\
Ronde 5  & $[235, 135, 145, 117, 12, 110, 48, 230, 232, 132, 252, 60, 13, 212, 180, 128]$ \\
Ronde 6  & $[138, 77, 173, 16, 75, 164, 46, 254, 98, 115, 159, 57, 162, 247, 138, 241]$ \\
Ronde 7  & $[168, 72, 140, 87, 85, 84, 195, 125, 159, 143, 106, 132, 226, 239, 39, 87]$ \\
Ronde 8  & $[243, 188, 166, 178, 129, 233, 188, 176, 201, 176, 13, 231, 222, 192, 167, 123]$ \\
Ronde 9  & $[240, 72, 254, 194, 17, 163, 132, 219, 195, 240, 136, 195, 208, 194, 145, 41]$ \\
\midrule
\textbf{Ronde 10} & $\mathbf{[83, 39, 8, 45, 30, 111, 246, 243, 213, 92, 171, 82, 45, 188, 92, 228]}$ \\
\bottomrule
\end{tabular}
\end{table}

\textbf{Teks Tersandi Akhir (Ciphertext $C$):}
\begin{itemize}[leftmargin=*]
    \item Nilai Desimal (16 Byte): $\mathbf{[83, 39, 8, 45, 30, 111, 246, 243, 213, 92, 171, 82, 45, 188, 92, 228]}$
    \item Biner (128 bit):\\
    \texttt{01010011 00100111 00001000 00101101 00011110 01101111 11110110 11110011}\\
    \texttt{11010101 01011100 10101011 01010010 00101101 10111100 01011100 11100100}
\end{itemize}

\newpage
% =========================================================================
\section{Mode Operasi Block Cipher dan Padding}
% =========================================================================

\subsection{Skema Padding PKCS\#7}
Jika panjang pesan $M$ bukan kelipatan ukuran blok $B$ (misal $B = 16$ byte untuk AES atau $B = 8$ byte untuk DES), sejumlah $p$ byte ditambahkan ke ujung pesan:
\begin{equation}
p = B - (|M| \bmod B)
\end{equation}
Setiap byte padding bernilai tepat $p$.
\begin{itemize}[leftmargin=*]
    \item Jika pesan kurang 3 byte dari batas blok: tambahkan 3 byte dengan nilai desimal 3 ($[3, 3, 3]$).
    \item \textbf{Kasus Khusus:} Jika panjang pesan sudah tepat kelipatan $B$, tetap harus ditambahkan satu blok penuh berisi $B$ byte bernilai $B$ (misal 16 byte bernilai desimal 16 untuk AES).
\end{itemize}

\subsection{Tiga Mode Operasi Utama}

\begin{figure}[H]
\centering
\begin{tcolorbox}[title=1. Electronic Codebook (ECB) -- \textbf{TIDAK AMAN}]
Setiap blok teks terang $P_i$ dienkripsi secara independen:
\begin{equation}
C_i = E_K(P_i), \quad P_i = D_K(C_i)
\end{equation}
\textbf{Kelemahan Fatal:} Dua blok teks terang yang sama akan selalu menghasilkan blok teks tersandi yang persis sama, sehingga membocorkan pola data.
\end{tcolorbox}
\end{figure}

\begin{figure}[H]
\centering
\begin{tcolorbox}[title=2. Cipher Block Chaining (CBC)]
Setiap blok teks terang di-XOR dengan blok ciphertext sebelumnya sebelum dienkripsi. Blok pertama di-XOR dengan Vektor Inisialisasi ($IV$) acak:
\begin{align}
C_0 &= IV \\
C_i &= E_K(P_i \oplus C_{i-1}), \quad i \ge 1 \\
P_i &= D_K(C_i) \oplus C_{i-1}, \quad i \ge 1
\end{align}
\textbf{Karakteristik:} Menghilangkan pola identik. $IV$ wajib acak dan tidak dapat diprediksi.
\end{tcolorbox}
\end{figure}

\begin{figure}[H]
\centering
\begin{tcolorbox}[title=3. Counter Mode (CTR)]
Mengubah block cipher menjadi stream cipher. Nilai counter yang unik dienkripsi, lalu hasilnya di-XOR dengan teks terang:
\begin{align}
T_i &= \text{Nonce} \parallel i \\
C_i &= P_i \oplus E_K(T_i) \\
P_i &= C_i \oplus E_K(T_i)
\end{align}
\textbf{Keunggulan:} Dapat diparalelkan 100\% dan fungsi dekripsi hanya membutuhkan fungsi enkripsi $E_K$.
\end{tcolorbox}
\end{figure}

\newpage
% =========================================================================
\section{Analisis Keamanan dan Efek Salju (Avalanche Effect)}
% =========================================================================

\subsection{Strict Avalanche Criterion (SAC)}
Jika 1 bit teks terang atau 1 bit kunci diubah, rata-rata **50\% bit teks tersandi harus berubah** secara acak.
\begin{itemize}[leftmargin=*]
    \item \textbf{Pada DES:} Membutuhkan sekitar 4 hingga 5 ronde untuk mencapai 50\% bit berubah.
    \item \textbf{Pada AES:} Berkat matriks MDS MixColumns, difusi 50\% bit flip tercapai hanya dalam \textbf{2 ronde saja}!
\end{itemize}

\subsection{Tabel Rangkuman Komparasi Teknis}
\begin{table}[H]
\centering
\caption{Matriks Perbandingan Mendalam DES vs AES}
\label{tab:deep_comparison}
\vspace{0.2cm}
\begin{tabular}{p{3.8cm}p{5.5cm}p{5.5cm}}
\toprule
\textbf{Aspek} & \textbf{Data Encryption Standard (DES)} & \textbf{Advanced Encryption Standard (AES)} \\
\midrule
\textbf{Ukuran Blok} & 64 bit (8 karakter teks) & 128 bit (16 karakter teks) \\
\textbf{Ukuran Kunci} & 56 bit efektif & 128, 192, 256 bit \\
\textbf{Fondasi Matematika} & Permutasi bit \& S-Box empiris & Aljabar Lapangan Terhingga $GF(2^8)$ \& Matriks MDS \\
\textbf{Struktur Inti} & Feistel (16 ronde) & SPN (10, 12, 14 ronde) \\
\textbf{Invertibilitas} & Fungsi ronde tidak harus bijektif & Setiap langkah harus memiliki fungsi invers \\
\textbf{Kekuatan Brute Force} & $2^{56}$ (rentan / dapat dibobol cepat) & $2^{128} \approx 3.4 \times 10^{38}$ (sangat aman) \\
\textbf{Efisiensi Software} & Lambat (operasi bit individu) & Sangat cepat (berbasis byte/word) \\
\bottomrule
\end{tabular}
\end{table}

\section{Panduan Praktikum \& Eksperimen Notebook (.ipynb)}
Dalam folder ini telah disediakan 3 buah Jupyter Notebook interaktif yang ditulis \textbf{100\% from scratch} menggunakan Python murni tanpa pustaka kriptografi eksternal:
\begin{enumerate}[leftmargin=*]
    \item \textbf{\texttt{01\_DES\_From\_Scratch.ipynb}:}
    Membangun manipulasi bit, tabel IP, PC-1, PC-2, ekspansi E, 8 Kotak-S DES, fungsi Feistel, 16 ronde enkripsi/dekripsi, dan uji efek salju.
    \item \textbf{\texttt{02\_AES\_From\_Scratch.ipynb}:}
    Membangun aritmetika Galois Field $GF(2^8)$, pembangkitan aljabar Kotak-S, transformasi $\mathtt{SubBytes}$, $\mathtt{ShiftRows}$, $\mathtt{MixColumns}$, $\mathtt{AddRoundKey}$, dan ekspansi kunci AES-128.
    \item \textbf{\texttt{03\_Mode\_Operasi\_dan\_Eksperimen\_Lanjutan.ipynb}:}
    Implementasi padding PKCS\#7, perbandingan pola ECB vs CBC vs CTR, dan eksperimen transmisi data rusak.
\end{enumerate}

\vspace{0.5cm}
\begin{center}
\textbf{-- Selamat Belajar dan Bereksperimen --}\\
\textit{"Kriptografi yang sejati bukan sekadar memanggil pustaka pustaka jadi, melainkan memahami keindahan matematika dan manipulasi bit di balik setiap lapisan transformasi."}
\end{center}

\end{document}
